\documentclass[10pt,a4paper]{amsart}
\usepackage[margin=19mm]{geometry}
\usepackage{amsmath,amssymb,mathtools}
\usepackage[scr=rsfs,cal=euler]{mathalfa}
\usepackage{xfrac}
\usepackage[unicode,hidelinks,bookmarks=false]{hyperref}
\newcommand{\ie}{i.e.\ }
\newcommand{\cf}{cf.\ }
\newcommand{\resp}{respectively\ }
\newcommand{\Subsection}{\subsection}
\providecommand{\nobreakdash}{\nobreak\hbox{-}}
\setlength{\emergencystretch}{2em}
\multlinegap0pt
\usepackage{graphics}
\usepackage{epsfig}
\usepackage{amscd}
\usepackage{xy}
\usepackage{latexsym}
\usepackage{multirow}
\newtheorem{thm}{Theorem}[section]
\newtheorem{lem}[thm]{Lemma}
\newtheorem{cor}[thm]{Corollary}
\newtheorem{prop}[thm]{Proposition}
\newtheorem{defn}[thm]{Definition}
\newtheorem{ex}[thm]{Example}
\newtheorem{es}{Exercise}
\newtheorem*{rem}{Remark}
\newtheorem*{pbm}{Problem}
\newtheorem*{Sol}{Sol.}
\newtheorem*{Q}{Question}
\newtheorem*{conje}{Conjecture}
\newtheorem{ass}[thm]{Assumption}
\newtheorem*{pf}{Proof}
\numberwithin{equation}{section} \numberwithin{figure}{section}
\renewcommand*{\to}{\rightarrow}
\renewcommand*{\bar}[1]{\overline{#1}}
\renewcommand{\parallel}{\mbox{ $\slash\!\slash$ }}
\renewcommand{\Re}{\operatorname{Re}}
\renewcommand{\Im}{\operatorname{Im}}
\newcommand{\Graph}{\operatorname{Graph}}
\newcommand{\Erf}{\operatorname{Erf}}
\newcommand{\Td}{\operatorname{Td}}
\newcommand{\res}{\operatorname{res}}
\newcommand{\Pf}{\operatorname{Pf}}
\newcommand{\gl}{\operatorname{gl}}
\newcommand{\codim}{\operatorname{codim}}
\newcommand{\alg}{\operatorname{alg}}
\newcommand{\norm}[1]{\left\|#1\right\|}
\newcommand{\Ric}{\operatorname{Ric}}
\newcommand{\ch}{\operatorname{ch}}
\newcommand{\Gal}{\operatorname{Gal}}
\newcommand{\sech}{\operatorname{sech}}
\newcommand{\csch}{\operatorname{csch}}
\newcommand{\vect}[1]{\widehat{\mathbf{#1}}}
\newcommand{\Diff}{\operatorname{Diff}}
\newcommand{\iso}{\operatorname{iso}}
\newcommand{\Aut}{\operatorname{Aut}}
\newcommand{\Aff}{\operatorname{Aff}}
\newcommand{\ind}{\operatorname{ind}}
\newcommand{\Gr}{\operatorname{Gr}}
\newcommand{\End}{\operatorname{End}}
\newcommand{\coker}{\operatorname{coker}}
\newcommand{\Res}{\operatorname{Res}}
\newcommand{\Vect}{\operatorname{Vect}}
\newcommand{\Hol}{\operatorname{Hol}}
\newcommand{\Har}{\operatorname{Har}}
\newcommand{\ord}{\operatorname{ord}}
\newcommand{\Gap}{\operatorname{Gap}}
\newcommand{\Sch}{\operatorname{Sch}}
\newcommand{\Sets}{\operatorname{Sets}}
\newcommand{\Mor}{\operatorname{Mor}}
\newcommand{\func}{\operatorname{func}}
\newcommand{\Herm}{\operatorname{Herm}}
\newcommand{\Ab}{\operatorname{Ab}}
\newcommand{\Ind}{\operatorname{Ind}}
\newcommand{\moduli}{\bar{\mc M}}
\newcommand{\uni}{\bar{\mc C}}
\newcommand{\vdim}{\operatorname{vdim}}
\newcommand{\virt}{\operatorname{virt}}
\newcommand{\ev}{\operatorname{ev}}
\newcommand{\st}[1]{\left|#1\right\rangle}
\newcommand{\dst}[1]{\left\langle #1\right|}
\newcommand{\hsp}{\left(\mc H,\ipd \cdot\cdot\right)}
\newcommand{\bsp}{\left(X,\|\cdot\|\right)}
\newcommand{\sgn}{\operatorname{sgn}}
\newcommand{\Sym}{\operatorname{Sym}}
\newcommand{\supp}{\operatorname{supp}}
\newcommand{\Lip}{\operatorname{Lip}}
\newcommand{\lebg}{(\mb R,\mc R, \alpha)}
\newcommand{\nlebg}{(\mb R^{n},\mc R^{n}, \alpha)}
\newcommand{\pbs}{(\Omega,\mc F,\mb P)}
\newcommand{\msp}{(\Omega,\mc F,\mu)}
\newcommand{\msk}[1]{(\Omega_{#1},\mc F_{#1},\mu_{#1})}
\newcommand{\var}{\operatorname{var}}
\newcommand{\dev}{\operatorname{dev}}
\renewcommand{\vec}[1]{\overrightarrow{#1}}
\newcommand{\tr}{\operatorname{Tr}}
\newcommand{\Sp}{\operatorname{span}}
\newcommand{\nullity}{\operatorname{nullity}}
\newcommand{\rank}{\operatorname{rank}}
\newcommand{\cof}{\operatorname{cof}}
\newcommand{\adj}{\operatorname{adj}}
\newcommand{\ipd}[2]{\langle #1|#2 \rangle} % inner product
\newcommand{\ran}{\operatorname{ran}}
\newcommand{\cod}{\operatorname{codim}}
\newcommand{\Ber}{\operatorname{Ber}}
\newcommand{\str}{\operatorname{str}}
\newcommand{\im}{\operatorname{Im}}
\newcommand{\ad}{\operatorname{ad}}
\newcommand{\der}{\operatorname{der}}
\newcommand{\mb}[1]{\mathbb{#1}} % Field
\newcommand{\stab}{\operatorname{stab}}
\newcommand{\Hom}{\operatorname{Hom}}
\newcommand{\hs}{\mathcal{H}}% Hilbert space
\newcommand{\bs}{\mathcal{X}}% Banach space
\newcommand{\bdo}[1]{\mathcal{B}(#1)}% The set of bounded operators
\newcommand{\fro}[1]{\mathcal{F}(#1)}% Finite rank operators
\newcommand{\mc}[1]{\mathcal{#1}}
\newcommand{\mk}[1]{\mathfrak{#1}}
\newcommand{\tsp}{(X,\tau)}
\newcommand{\Det}{\operatorname{Det}}
\newcommand{\Tr}{\operatorname{Tr}}
\newcommand{\lk}{\operatorname{link}}
\newcommand{\Der}{\operatorname{Der}}
\newcommand{\Spec}{\operatorname{Spec}}
\newcommand{\Tor}{\operatorname{Tor}}
\newcommand{\Coh}{\operatorname{Coh}}
\newcommand{\Ht}{\operatorname{Ht}}
\newcommand{\id}{\operatorname{id}}
\newcommand{\ob}{\operatorname{ob}}
\newcommand{\hto}{\hookrightarrow}
\newcommand{\tto}{\twoheadrightarrow}
\newcommand{\Pic}{\operatorname{Pic}}
\newcommand{\St}{\operatorname{St}}
\newcommand{\Ext}{\operatorname{Ext}}
\newcommand{\Sh}{\operatorname{Sh}}
\newcommand{\diag}{\operatorname{diag}}
\newcommand{\Alt}{\operatorname{Alt}}
\newcommand{\dbar}{\bar{\partial}}
\renewcommand{\d}{\partial}
\newcommand{\dsp}{\displaystyle}
\newcommand{\mf}[1]{\mathbf{#1}}
\newcommand{\Jac}{\operatorname{Jac}}
\usepackage{url}
\usepackage{pgfplots}
\pgfplotsset{compat=1.18}
\usetikzlibrary{calc} 
\usepackage{mathtools}
\usepackage{tikz-cd}
\begin{document}
\title[Weyl Curves and Zeta Determinants of Conic Laplacians on Rie]{Weyl Curves and Zeta Determinants of Conic Laplacians on Riemann Surfaces}
\author{Jia-Ming (Frank) Liou}
\date{}
\maketitle
\noindent\emph{Source and licence.} Weyl Curves and Zeta Determinants of Conic Laplacians on Riemann Surfaces. Version arXiv:2606.20818v2 (CC BY 4.0). This material is distributed under the Creative Commons Attribution 4.0 International licence, http://creativecommons.org/licenses/by/4.0/, as recorded in the arXiv OAI licence metadata for this exact version and shown on the abstract page https://arxiv.org/abs/2606.20818v2. Full rights notice: licenses/arxiv-2606.20818-CC-BY-4.0.txt.\par\smallskip
\noindent\textit{Editorial scope.} Counted unit: the original \texttt{thm} environment \texttt{thm:weyl-transformation-law} at source lines 3835--3873 together with its complete proof at lines 3875--3948; the delivered document keeps source lines 3192--3949 so that every referenced label and every citation resolves inside it. Native editable source is the paper's own concatenated TeX; the AMS wrapper is editorial formatting only. No publisher class, upstream script or unrelated artwork is redistributed.
\section{The Weyl Curve and the Determinant Line}
We now organize the preceding analytic results geometrically. Related
maps into Grassmannians, formed from null spaces of elliptic cone
operators, appear in \cite{GilKrainerMendoza2007}, while Wang developed an abstract
theory of Weyl curves for symmetric operators
\cite{WangWeylCurve}. The treatment below is self-contained and is
formulated directly in terms of the critical asymptotic boundary data
constructed above. Its purpose is to interpret the analytic identities
from the preceding sections: local graph coordinates recover the Weyl
functions, the real tangent form recovers their derivative identity,
and the pullback determinant line gives a geometric interpretation of
the boundary and zeta determinants.

For \(\lambda\in\mb C\), define the intrinsic boundary-data
subspace
\(\mc K_X(\lambda):=
\pi_S\bigl(\ker(\Delta_X-\lambda)\bigr)\subset V_S\).

\begin{lem}\label{lem:cauchy-data-orthogonality}
For every \(\lambda\in\mb C\), one has
\(\mc K_X(\overline{\lambda})
\subset\mc K_X(\lambda)^{\perp_{\omega_S}}\).
In particular, if \(\lambda\in\mb R\), then
\(\mc K_X(\lambda)\) is isotropic.
\end{lem}

\begin{proof}
Let \(g\in\mc K_X(\overline{\lambda})\) and
\(h\in\mc K_X(\lambda)\). Choose
\(v\in\ker(\Delta_X-\overline{\lambda})\) and
\(u\in\ker(\Delta_X-\lambda)\) such that
\(\pi_Sv=g\) and \(\pi_Su=h\). By the Green identity,
\[
\begin{aligned}
\omega_S(g,h)
&=
q_X(v,u) \\
&=
\langle\Delta_Xv,u\rangle
-
\langle v,\Delta_Xu\rangle \\
&=
\overline{\lambda}\langle v,u\rangle
-
\overline{\lambda}\langle v,u\rangle
=
0.
\end{aligned}
\]
Here the second equality in the last line uses the convention that
the inner product is linear in the first variable and conjugate
linear in the second. Hence
\(g\in\mc K_X(\lambda)^{\perp_{\omega_S}}\), proving the asserted
inclusion.

If \(\lambda\in\mb R\), then
\(\overline{\lambda}=\lambda\), and hence
\(\mc K_X(\lambda)
\subset\mc K_X(\lambda)^{\perp_{\omega_S}}\).
Thus \(\mc K_X(\lambda)\) is isotropic.
\end{proof}

\begin{lem}\label{lem:cauchy-data-orthogonal-complement}
Let \(\mc V\in\operatorname{LGr}(V_S,\omega_S)\). If
\(\lambda\in\rho(\Delta_{X,\mc V})\), then
\(\mc K_X(\overline{\lambda})
=\mc K_X(\lambda)^{\perp_{\omega_S}}\).
In particular, if
\(\lambda\in\mb R\cap\rho(\Delta_{X,\mc V})\), then
\(\mc K_X(\lambda)\) is Lagrangian.
\end{lem}
\begin{proof}
By Lemma~\ref{lem:cauchy-data-orthogonality},
\(\mc K_X(\overline{\lambda})
\subset\mc K_X(\lambda)^{\perp_{\omega_S}}\).
Choose a Lagrangian complement \(\mc V^\#\) of \(\mc V\). Since
\(\Delta_{X,\mc V}\) is self-adjoint,
\(\lambda\in\rho(\Delta_{X,\mc V})\) implies
\(\overline{\lambda}\in\rho(\Delta_{X,\mc V})\). For
\(\mu\in\{\lambda,\overline{\lambda}\}\), the Poisson operator
\(\Gamma_{\mc V,\mc V^\#}^X(\mu):
\mc V^\#\to\ker(\Delta_X-\mu)\)
is an isomorphism.
Moreover, the restriction of \(\pi_S\) to
\(\ker(\Delta_X-\mu)\) is injective. Indeed, if
\(u\in\ker(\Delta_X-\mu)\) and \(\pi_Su=0\), then
\(u\in D_{\min}(\Delta_X)\subset D_{\mc V}\), and hence
\((\Delta_{X,\mc V}-\mu)u=0\). Since
\(\mu\in\rho(\Delta_{X,\mc V})\), it follows that \(u=0\).
By the definition of \(\mc K_X(\mu)\), this restriction is also
surjective onto \(\mc K_X(\mu)\), and is therefore an isomorphism.
Consequently,
\[
\dim_{\mb C}\mc K_X(\lambda)
=
\dim_{\mb C}\mc K_X(\overline{\lambda})
=
\dim_{\mb C}\mc V^\#
=
\frac12\dim_{\mb C}V_S.
\]
Since \(\omega_S\) is nondegenerate,
\(\dim_{\mb C}\mc K_X(\lambda)^{\perp_{\omega_S}}
=\dim_{\mb C}V_S-\dim_{\mb C}\mc K_X(\lambda)
=\frac12\dim_{\mb C}V_S\).
The preceding inclusion is therefore an equality.

If \(\lambda\in\mb R\), then
\(\overline{\lambda}=\lambda\), and hence
\(\mc K_X(\lambda)
=\mc K_X(\lambda)^{\perp_{\omega_S}}\).
Thus \(\mc K_X(\lambda)\) is Lagrangian.
\end{proof}

Set
\(
    \Sigma_X
    :=
    \bigcup_{\mc V\in\operatorname{LGr}(V_S,\omega_S)}
    \rho(\Delta_{X,\mc V}).
\)

\begin{prop}\label{prop:domain-of-cauchy-data-map}
We have
\[
    \Sigma_X
    =
    \mb C\setminus\sigma_p(\Delta_{X,\min}),
\]
where \(\sigma_p(\Delta_{X,\min})\) denotes the point spectrum of the
minimal Laplacian. Consequently, \(\Sigma_X\) is open and
path-connected.
\end{prop}

\begin{proof}
Since
\(
    \Sigma_X
    =
    \bigcup_{\mc V\in\operatorname{LGr}(V_S,\omega_S)}
    \rho(\Delta_{X,\mc V}),
\)
and each resolvent set is open, the set \(\Sigma_X\) is open.
Suppose first that
\(\lambda_0\in\sigma_p(\Delta_{X,\min})\). Then there exists
\(0\neq u_0\in D_{\min}(\Delta_X)\) such that
\(
    (\Delta_{X,\min}-\lambda_0)u_0=0.
\)
For every
\(\mc V\in\operatorname{LGr}(V_S,\omega_S)\), we have
\(D_{\min}(\Delta_X)\subset D_{\mc V}\), and hence
\(
    (\Delta_{X,\mc V}-\lambda_0)u_0=0.
\)
Thus
\(\lambda_0\in\sigma_p(\Delta_{X,\mc V})\) for every \(\mc V\), so
\(\lambda_0\notin\Sigma_X\). Therefore
\[
    \sigma_p(\Delta_{X,\min})
    \subset
    \mb C\setminus\Sigma_X.
\]
Conversely, suppose that
\(\lambda_0\notin\sigma_p(\Delta_{X,\min})\). We show that
\(\lambda_0\in\Sigma_X\).
If \(\lambda_0\notin\mb R\), then
\(\lambda_0\in\rho(\Delta_{X,\mc V})\) for every
\(\mc V\in\operatorname{LGr}(V_S,\omega_S)\), since each
\(\Delta_{X,\mc V}\) is self-adjoint. Hence
\(\lambda_0\in\Sigma_X\).
Suppose now that \(\lambda_0\in\mb R\). Then
\(\mc K_X(\lambda_0)\) is isotropic. Choose a Lagrangian subspace
\(\mc W\subset V_S\) containing \(\mc K_X(\lambda_0)\), and choose a
Lagrangian complement \(\mc W^\#\) of \(\mc W\). It follows that
\(
    \mc K_X(\lambda_0)\cap\mc W^\#=\{0\}.
\)
We claim that \(\lambda_0\) is not an eigenvalue of
\(\Delta_{X,\mc W^\#}\). Indeed, suppose that
\(u\in D_{\mc W^\#}\) satisfies
\(
    (\Delta_{X,\mc W^\#}-\lambda_0)u=0.
\)
Then
\(
    \pi_Su
    \in
    \mc K_X(\lambda_0)\cap\mc W^\#
    =
    \{0\}.
\)
Hence \(u\in D_{\min}(\Delta_X)\) and
\(
    (\Delta_{X,\min}-\lambda_0)u=0.
\)
Since
\(\lambda_0\notin\sigma_p(\Delta_{X,\min})\), we conclude that
\(u=0\). Thus
\(\lambda_0\notin\sigma_p(\Delta_{X,\mc W^\#})\).
Since \(\Delta_{X,\mc W^\#}\) is self-adjoint with compact resolvent,
its spectrum consists entirely of isolated eigenvalues of finite
multiplicity. Therefore
\(\lambda_0\in\rho(\Delta_{X,\mc W^\#})\), and hence
\(\lambda_0\in\Sigma_X\). This proves
\[
    \Sigma_X
    =
    \mb C\setminus\sigma_p(\Delta_{X,\min}).
\]
Finally, since \(\Delta_{X,\min}\) is symmetric,
\(\sigma_p(\Delta_{X,\min})\subset\mb R\). Moreover, for any fixed
\(\mc V\in\operatorname{LGr}(V_S,\omega_S)\),
\[
    \sigma_p(\Delta_{X,\min})
    \subset
    \sigma(\Delta_{X,\mc V}).
\]
Since \(\Delta_{X,\mc V}\) has compact resolvent, its spectrum has no
finite accumulation point. Consequently, the point spectrum of
\(\Delta_{X,\min}\) is a closed discrete subset of \(\mb R\). Its
complement in \(\mb C\) is path-connected.
Therefore \(\Sigma_X\) is path-connected.
\end{proof}
Set \(m:=\frac12\dim_{\mb C}V_S\), and let
\(\operatorname{Gr}_m(V_S)\) be the complex Grassmannian. Its points
are the \(m\)-dimensional complex linear subspaces of \(V_S\), and its
complex dimension is \(m^2\).

We regard the Lagrangian Grassmannian
\(\operatorname{LGr}(V_S,\omega_S)\) as a subset of
\(\operatorname{Gr}_m(V_S)\).

For every \(\lambda\in\Sigma_X\), there exists
\(\mc V\in\operatorname{LGr}(V_S,\omega_S)\) such that
\(\lambda\in\rho(\Delta_{X,\mc V})\). Choose a Lagrangian complement
\(\mc V^\#\) of \(\mc V\). The Poisson operator is an isomorphism
from \(\mc V^\#\) onto \(\ker(\Delta_X-\lambda)\). Hence
\(\dim_{\mb C}\ker(\Delta_X-\lambda)=m\).
Moreover, the restriction
\[
    \pi_S\big|_{\ker(\Delta_X-\lambda)}:
    \ker(\Delta_X-\lambda)
    \longrightarrow
    \mc K_X(\lambda)
\]
is a linear isomorphism. Indeed, it is surjective by the definition of
\(\mc K_X(\lambda)\). If
\(u\in\ker(\Delta_X-\lambda)\) and \(\pi_Su=0\), then
\(u\in D_{\min}(\Delta_X)\subset D_{\mc V}\), and hence
\(
    (\Delta_{X,\mc V}-\lambda)u=0.
\)
Since \(\lambda\in\rho(\Delta_{X,\mc V})\), it follows that \(u=0\).
Therefore
\[
    \dim_{\mb C}\mc K_X(\lambda)
    =
    \dim_{\mb C}\ker(\Delta_X-\lambda)
    =
    m.
\]
We thus obtain a map
\[
    \mc K_X:
    \Sigma_X
    \longrightarrow
    \operatorname{Gr}_m(V_S),
    \qquad
    \lambda
    \longmapsto
    \mc K_X(\lambda).
\]
We next prove that \(\mc K_X\) is holomorphic. We first establish the
local graph representation needed for this purpose.


Let \(E\in\operatorname{Gr}_m(V_S)\), and choose
\(F\in\operatorname{Gr}_m(V_S)\) such that
\(V_S=E\oplus F\). Consider the open subset
\[
    \mc U_E
    =
    \left\{
        \mc L\in\operatorname{Gr}_m(V_S):
        \mc L\cap E=\{0\}
    \right\}.
\]
Let \(P_E:V_S\to E\) and \(P_F:V_S\to F\) be the projections
associated with this decomposition.
For every \(\mc L\in\mc U_E\), the restriction
\(P_F|_{\mc L}:\mc L\to F\) is an isomorphism. Define
\(
    A_{\mc L}
    :=
    P_E\circ
    \bigl(P_F|_{\mc L}\bigr)^{-1}
    \in
    \operatorname{Hom}_{\mb C}(F,E).
\)
Then, for every \(h\in F\),
\(
    \bigl(P_F|_{\mc L}\bigr)^{-1}h
    =
    A_{\mc L}h+h,
\)
and hence
\[
    \mc L
    =
    \left\{
        A_{\mc L}h+h:
        h\in F
    \right\}
    =
    \operatorname{graph}_{(E,F)}(A_{\mc L}).
\]
The correspondence
\[
    \kappa_{(E,F)}:
    \mc U_E
    \longrightarrow
    \operatorname{Hom}_{\mb C}(F,E),
    \qquad
    \mc L
    \longmapsto
    A_{\mc L},
\]
defines a coordinate chart on
\(\operatorname{Gr}_m(V_S)\).

\begin{lem}\label{lem:cauchy-data-transverse}
Let \(\mc V\in\operatorname{LGr}(V_S,\omega_S)\). For every
\(\lambda\in\rho(\Delta_{X,\mc V})\),
\[
    \mc K_X(\lambda)\cap\mc V=\{0\}.
\]
\end{lem}

\begin{proof}
Let \(h\in\mc K_X(\lambda)\cap\mc V\). By the definition of
\(\mc K_X(\lambda)\), there exists
\(u_\lambda\in\ker(\Delta_X-\lambda)\) such that
\(\pi_Su_\lambda=h\). Since \(h\in\mc V\), we have
\(u_\lambda\in D_{\mc V}\). Hence
\(
    (\Delta_{X,\mc V}-\lambda)u_\lambda=0.
\)
Since \(\lambda\in\rho(\Delta_{X,\mc V})\),
the operator \(\Delta_{X,\mc V}-\lambda\) is invertible. Therefore
\(u_\lambda=0\), and consequently \(h=\pi_Su_\lambda=0\).
\end{proof}



Let \((\mc V,\mc V^\#)\) be a Lagrangian pair. If
\(\lambda\in\rho(\Delta_{X,\mc V})\), then
\(\mc K_X(\lambda)\cap\mc V=\{0\}\), and hence
\(\mc K_X(\lambda)\in\mc U_{\mc V}\). By the definition of the Weyl
function,
\[
    \mc K_X(\lambda)
    =
    \operatorname{graph}_{(\mc V,\mc V^\#)}
    \bigl(M_{\mc V,\mc V^\#}^X(\lambda)\bigr).
\]
Consequently,
\[
    \kappa_{(\mc V,\mc V^\#)}
    \bigl(\mc K_X(\lambda)\bigr)
    =
    M_{\mc V,\mc V^\#}^X(\lambda)
    \in
    \operatorname{Hom}_{\mb C}(\mc V^\#,\mc V).
\]

\begin{prop}
The map
\(
    \mc K_X:
    \Sigma_X
    \longrightarrow
    \operatorname{Gr}_m(V_S)
\)
is holomorphic.
\end{prop}

\begin{proof}
Let \(\lambda_0\in\Sigma_X\). By the definition of \(\Sigma_X\), there
exists \(\mc V\in\operatorname{LGr}(V_S,\omega_S)\) such that
\(\lambda_0\in\rho(\Delta_{X,\mc V})\). Choose a Lagrangian complement
\(\mc V^\#\) of \(\mc V\).

For every \(\lambda\in\rho(\Delta_{X,\mc V})\), the previous lemma
gives \(\mc K_X(\lambda)\cap\mc V=\{0\}\). Hence
\(
    \mc K_X\bigl(\rho(\Delta_{X,\mc V})\bigr)
    \subset
    \mc U_{\mc V}.
\)
Consider the identity chart on
\(\rho(\Delta_{X,\mc V})\subset\mb C\) and the Grassmannian chart
\(
    \kappa_{(\mc V,\mc V^\#)}:
    \mc U_{\mc V}
    \longrightarrow
    \operatorname{Hom}_{\mb C}(\mc V^\#,\mc V).
\)
The corresponding local coordinate representation of \(\mc K_X\) is
\[
    \kappa_{(\mc V,\mc V^\#)}
    \circ
    \mc K_X\big|_{\rho(\Delta_{X,\mc V})}
    =
    M_{\mc V,\mc V^\#}^X.
\]
Since \(M_{\mc V,\mc V^\#}^X\) is holomorphic on
\(\rho(\Delta_{X,\mc V})\), the map \(\mc K_X\) is holomorphic at
\(\lambda_0\). Since \(\lambda_0\) was arbitrary, \(\mc K_X\) is
holomorphic on \(\Sigma_X\).
\end{proof}

\begin{defn}
The holomorphic map
\[
    \mc K_X:
    \Sigma_X
    \longrightarrow
    \operatorname{Gr}_m(V_S)
\]
is called the Weyl curve of \(X\).
\end{defn}

Differentiating the local coordinate identity established above
immediately gives the following description of the differential of
the Weyl curve.

\begin{prop}\label{prop:differential-weyl-curve}
Let \(\lambda_0\in\Sigma_X\). The holomorphic differential of the
Weyl curve at \(\lambda_0\) is the complex-linear map
\[
    d\mc K_X|_{\lambda_0}:
    T_{\lambda_0}^{1,0}\Sigma_X
    \longrightarrow
    T_{\mc K_X(\lambda_0)}^{1,0}
    \operatorname{Gr}_m(V_S).
\]
Let \((\mc V,\mc V^\#)\) be a Lagrangian pair such that
\(\lambda_0\in\rho(\Delta_{X,\mc V})\). Under the natural
identification
\[
    T_{\mc K_X(\lambda_0)}^{1,0}
    \operatorname{Gr}_m(V_S)
    \simeq
    \operatorname{Hom}_{\mb C}
    \bigl(
        \mc K_X(\lambda_0),
        V_S/\mc K_X(\lambda_0)
    \bigr),
\]
the tangent vector
\(
    d\mc K_X|_{\lambda_0}
    \bigl(
        \left.\frac{\partial}{\partial\lambda}\right|_{\lambda_0}
    \bigr)
\)
is given by
\[
\begin{aligned}
    &d\mc K_X|_{\lambda_0}
    \left(
        \left.\frac{\partial}{\partial\lambda}\right|_{\lambda_0}
    \right)
    \left(
        M_{\mc V,\mc V^\#}^X(\lambda_0)h^\#+h^\#
    \right)
    \\
    &\qquad
    =
    \left[
        \left.
        \frac{d}{d\lambda}
        M_{\mc V,\mc V^\#}^X(\lambda)h^\#
        \right|_{\lambda=\lambda_0}
    \right]
\end{aligned}
\]
for every \(h^\#\in\mc V^\#\), where the brackets denote the class in
\(V_S/\mc K_X(\lambda_0)\).

Equivalently, in the graph-coordinate chart determined by
\((\mc V,\mc V^\#)\),
\[
    d\kappa_{(\mc V,\mc V^\#)}|_{\mc K_X(\lambda_0)}
    \left(
        d\mc K_X|_{\lambda_0}
        \left(
            \left.\frac{\partial}{\partial\lambda}\right|_{\lambda_0}
        \right)
    \right)
    =
    \left.
    \frac{d}{d\lambda}
    M_{\mc V,\mc V^\#}^X(\lambda)
    \right|_{\lambda=\lambda_0}.
\]
\end{prop}

\begin{proof}
Fix \(h^\#\in\mc V^\#\). Near \(\lambda_0\), the vector
\(
    s_{h^\#}(\lambda)
    =
    M_{\mc V,\mc V^\#}^X(\lambda)h^\#+h^\#
\)
belongs to \(\mc K_X(\lambda)\). Differentiating at \(\lambda_0\)
gives
\[
    s_{h^\#}'(\lambda_0)
    =
    \left.
    \frac{d}{d\lambda}
    M_{\mc V,\mc V^\#}^X(\lambda)h^\#
    \right|_{\lambda=\lambda_0}.
\]
Under the standard identification of the holomorphic tangent space of
the Grassmannian with
\(
    \operatorname{Hom}_{\mb C}
    \bigl(
        \mc K_X(\lambda_0),
        V_S/\mc K_X(\lambda_0)
    \bigr),
\)
the corresponding tangent vector is obtained by taking the class of
\(s_{h^\#}'(\lambda_0)\) modulo \(\mc K_X(\lambda_0)\). This proves
the first formula. The coordinate formula follows by differentiating
\(
    \kappa_{(\mc V,\mc V^\#)}\circ\mc K_X
    =
    M_{\mc V,\mc V^\#}^X
\)
at \(\lambda_0\).
\end{proof}

We now study the transition between the local coordinate
representations of the Weyl curve determined by two Lagrangian pairs
\((\mc V,\mc V^\#)\) and \((\mc W,\mc W^\#)\) in
\((V_S,\omega_S)\).

Let \(V,W,E,F\) be subspaces of \(V_S\) such that
\(
    V_S=V\oplus W=E\oplus F.
\)
Every linear map \(T:V_S\to V_S\) admits a block representation
relative to the source decomposition \(V_S=V\oplus W\) and the target
decomposition \(V_S=E\oplus F\). Namely, there exist linear maps
\(
    T_{11}:V\to E,
    T_{12}:W\to E,
    T_{21}:V\to F,
\)
and \(T_{22}:W\to F\) such that, for \(v\in V\) and \(w\in W\),
\[
    T(v+w)
    =
    \bigl(T_{11}v+T_{12}w\bigr)
    +
    \bigl(T_{21}v+T_{22}w\bigr),
\]
where the first term belongs to \(E\) and the second belongs to \(F\).
We denote this block representation by
\[
    [T]_{(V,W)}^{(E,F)}
    =
    \begin{bmatrix}
        T_{11} & T_{12}\\
        T_{21} & T_{22}
    \end{bmatrix}.
\]

Assume that
\(
    \dim_{\mb C}V
    =
    \dim_{\mb C}W
    =
    \dim_{\mb C}E
    =
    \dim_{\mb C}F
    =
    m.
\)
Consider the transition map
\(
    \kappa_{(E,F)}
    \circ
    \kappa_{(V,W)}^{-1}
\)
between the graph-coordinate charts determined by the decompositions
\(V_S=V\oplus W\) and \(V_S=E\oplus F\).
Its domain is the open subset of
\(\operatorname{Hom}_{\mb C}(W,V)\) corresponding to
\(\mc U_V\cap\mc U_E\), and its values lie in
\(\operatorname{Hom}_{\mb C}(F,E)\).

Write
\[
    [\id_{V_S}]_{(V,W)}^{(E,F)}
    =
    \begin{bmatrix}
        A&B\\
        C&D
    \end{bmatrix}.
\]
The standard transition formula for graph coordinates is
\[
    \bigl(
        \kappa_{(E,F)}
        \circ
        \kappa_{(V,W)}^{-1}
    \bigr)(L)
    =
    (AL+B)(CL+D)^{-1}.
\]
Moreover,
\[
    \kappa_{(V,W)}
    \bigl(\mc U_V\cap\mc U_E\bigr)
    =
    \left\{
        L\in\operatorname{Hom}_{\mb C}(W,V):
        CL+D:W\longrightarrow F
        \text{ is invertible}
    \right\}.
\]



We therefore obtain the following transformation law for the local
coordinate representations of the Weyl curve.


\begin{thm}\label{thm:weyl-transformation-law}
Let \((\mc V,\mc V^\#)\) and \((\mc W,\mc W^\#)\) be Lagrangian
pairs. Let
\[
    [\id_{V_S}]_{(\mc V,\mc V^\#)}^{(\mc W,\mc W^\#)}
    =
    \begin{bmatrix}
        A&B\\
        C&D
    \end{bmatrix}
\]
be the block representation of the identity map relative to the two
Lagrangian decompositions. If
\(
    \lambda
    \in
    \rho(\Delta_{X,\mc V})
    \cap
    \rho(\Delta_{X,\mc W}),
\)
then
\(
    C M_{\mc V,\mc V^\#}^X(\lambda)+D:
    \mc V^\#
    \longrightarrow
    \mc W^\#
\)
is invertible, and
\[
    M_{\mc W,\mc W^\#}^X(\lambda)
    =
    \bigl(
        A M_{\mc V,\mc V^\#}^X(\lambda)+B
    \bigr)
    \bigl(
        C M_{\mc V,\mc V^\#}^X(\lambda)+D
    \bigr)^{-1}.
\]
\end{thm}

\begin{proof}
We argue directly from the definition of the Weyl function.

Let \(h^\#\in\mc V^\#\), and set
\(
    u_\lambda
    =
    \Gamma_{\mc V,\mc V^\#}^X(\lambda)h^\#.
\)
Then
\(\widetilde\pi_{\mc V^\#}(u_\lambda)=h^\#\), and
\[
    \pi_S(u_\lambda)
    =
    M_{\mc V,\mc V^\#}^X(\lambda)h^\#+h^\#.
\]
With respect to the decomposition
\(V_S=\mc W\oplus\mc W^\#\), its boundary components are
\[
\begin{aligned}
    \widetilde\pi_{\mc W}(u_\lambda)
    &=
    \bigl(
        A M_{\mc V,\mc V^\#}^X(\lambda)+B
    \bigr)h^\#,
    \\
    \widetilde\pi_{\mc W^\#}(u_\lambda)
    &=
    \bigl(
        C M_{\mc V,\mc V^\#}^X(\lambda)+D
    \bigr)h^\#.
\end{aligned}
\]
Set
\(
    k^\#
    =
    \bigl(
        C M_{\mc V,\mc V^\#}^X(\lambda)+D
    \bigr)h^\#.
\)
Since \(\lambda\in\rho(\Delta_{X,\mc W})\), the defining property of
the Poisson operator implies that
\(
    u_\lambda
    =
    \Gamma_{\mc W,\mc W^\#}^X(\lambda)k^\#.
\)

If \(k^\#=0\), then \(u_\lambda=0\), and hence
\(h^\#=\widetilde\pi_{\mc V^\#}(u_\lambda)=0\). Therefore
\(C M_{\mc V,\mc V^\#}^X(\lambda)+D\) is injective. Since
\(\dim_{\mb C}\mc V^\#=\dim_{\mb C}\mc W^\#=m\), it is invertible.

By the definition of the Weyl function,
\[
    M_{\mc W,\mc W^\#}^X(\lambda)k^\#
    =
    \bigl(
        A M_{\mc V,\mc V^\#}^X(\lambda)+B
    \bigr)h^\#.
\]
Since this holds for every \(h^\#\in\mc V^\#\), we obtain
\[
    M_{\mc W,\mc W^\#}^X(\lambda)
    \bigl(
        C M_{\mc V,\mc V^\#}^X(\lambda)+D
    \bigr)
    =
    A M_{\mc V,\mc V^\#}^X(\lambda)+B.
\]
The asserted transformation formula now follows from the
invertibility established above.
\end{proof}


\begin{thebibliography}{99}
\bibitem[GilKrainerMendoza2007]{GilKrainerMendoza2007} J. B. Gil, T. Krainer and G. A. Mendoza, Geometry and spectra of closed extensions of elliptic cone operators, Canad. J. Math. 59 (2007), no. 4, 742--794.
\bibitem[WangWeylCurve]{WangWeylCurve} Yicao Wang, Complex analysis of symmetric operators, I, arXiv:2408.10968v2 [math.FA], 2024.
\end{thebibliography}
\end{document}
